\section{Results}
\label{sec:results}

We evaluate where membership evidence remains after the SAE decomposition. The results build in three steps. First, membership-related knowledge
directions and recall-related directions are weakly aligned at the SAE layer, so membership detection and extractable recall are not the same geometric
object. Second, SAE reconstruction weakens membership detection, while detectors recover much of the lost signal from the reconstruction residual. Third,
feature-level and direction-targeted interventions can suppress some visible outputs without removing residual membership evidence.

\subsection{Membership and recall directions are weakly aligned}
\label{sec:results:geometry}

Across the nine fine-tuned transformers, the leading knowledge direction $\dK$ and the leading recall direction $\dR$ are weakly aligned at the sparse
autoencoder layer (Table~\ref{tab:bcd_main}, Appendix~\ref{app:per_model_bcd}). Principal angles between the corresponding subspaces are large across
models, and independently trained runs indicate that the cosine and angle estimates are stable. The main implication is that membership detection and
extractable recall do not reduce to the same leading activation direction.

Layer-by-layer detector curves support the same separation dynamically. Membership and recall signals develop on different schedules, with mutual alignment
rising in late layers only for the two models under capacity pressure (Figure~\ref{fig:layer_trajectories}, Appendix~\ref{app:per_model_bcd}). This
motivates the SAE reconstruction test. If the learned SAE representation preserves the membership-relevant part of the activation, reconstruction should
preserve membership detection performance. If reconstruction weakens detection while the residual preserves it, then feature-level interventions may miss
membership evidence left outside the reconstruction.

\subsection{SAE reconstruction fails to preserve membership signal recoverable from the residual}
\label{sec:results:dark}

We next test whether membership signal is preserved in the SAE reconstruction or remains recoverable from the reconstruction residual. Table~\ref{tab:dark_subspace} compares membership detection on three representations of the same activation, namely the original hidden state $\mathbf{h}$, the
SAE reconstruction $\hat{\mathbf{h}}$, and the reconstruction residual $\mathbf{h}-\hat{\mathbf{h}}$.

\begin{table*}[!t]
\centering
\caption{
Membership signal under SAE decomposition.
AUROC is reported for membership detectors trained on the original activation, the SAE reconstruction, and the reconstruction residual.
$\Delta_{\mathrm{recon}}$ is the change from Original to SAE-Recon.
Cosine reconstruction quality is reported as $\cos(x,\hat{x})$.
Settings that pass the validity gate (no $\dagger$ or $\ddagger$) form the set scored by the primary tests. The $\ddagger$-marked settings fall below the strict reconstruction-cosine threshold and are retained for diagnostic context, not as clean decomposition evidence.
Bootstrap confidence intervals and additional reconstruction-quality diagnostics are reported in Appendix~\ref{app:koc2_bootstrap}.
}
\label{tab:dark_subspace}
\small
\begin{tabular}{@{}lccccc@{}}
\toprule
\textbf{Model / SAE setting}
& \textbf{Original}
& \textbf{SAE-Recon}
& \textbf{Residual}
& $\boldsymbol{\Delta_{\mathrm{recon}}}$
& $\boldsymbol{\cos(x,\hat{x})}$ \\
\midrule
\multicolumn{6}{@{}l}{\emph{Controlled Pythia experiments}} \\
Pythia-6.9B                    & 0.803 & 0.593 & 0.805 & $-$21 pp & 0.927 \\
Pythia-6.9B (mixed)            & 0.803 & 0.594 & 0.781 & $-$21 pp & 0.976 \\
Pythia-12B                     & 0.764 & 0.612 & 0.723 & $-$15 pp & 0.991 \\
\midrule
\multicolumn{6}{@{}l}{\emph{Cross-architecture sweep}} \\
Pythia-1B$^{\ddagger}$         & 0.663 & 0.593 & 0.670 & $-$7 pp & 0.869 \\
GPT-Neo-2.7B                   & 0.615 & 0.571 & 0.699 & $-$4 pp  & 0.998 \\
GPT-Neo-2.7B (mixed)           & 0.616 & 0.558 & 0.646 & $-$6 pp  & 0.997 \\
OPT-6.7B$^{\ddagger}$          & 0.774 & 0.557 & 0.866$^{*}$ & $-$22 pp & 0.533 \\
Qwen2-7B                       & 0.638 & 0.526 & 0.849$^{*}$ & $-$11 pp & 0.850 \\
Falcon-7B$^{\dagger\ddagger}$ & 0.690 & 0.694 & 0.827$^{*}$ & $+$0.4 pp & 0.767 \\
Gemma-2-2B                     & 0.801 & 0.769 & 0.820 & $-$3 pp  & 0.911 \\
\bottomrule
\end{tabular}

\vspace{0.5em}
\footnotesize{
$^{*}$Residual norm AUROC exceeds 0.65.
$^{\dagger}$Falcon-7B activates only 9 SAE features.
$^{\ddagger}$Reconstruction cosine falls below the 0.85 permissive threshold, so the setting is retained as a diagnostic rather than as evidence eligible for quantitative comparison.
Token-level Bag-of-Words controls are reported in Appendix~\ref{app:bow_baseline}.
}
\end{table*}

The controlled Pythia experiments provide the strongest decomposition evidence. On Pythia-6.9B, both the member-only SAE and the label-balanced mixed-data
SAE weaken membership detection by 21 percentage points after reconstruction, while residual detectors recover much of the lost signal. Because the mixed-data SAE is not trained only on member activations, the matching reconstruction-induced reduction rules out member-only dictionary training as the sole explanation in this controlled setting. Pythia-12B shows the same ordering at smaller magnitude. The Pythia results therefore suggest that SAE
reconstruction discards membership evidence that remains available in the reconstruction residual.

The broader sweep tests whether this ordering is specific to Pythia. Among model-SAE settings where reconstruction is reliable enough to interpret, the same qualitative pattern appears in GPT-Neo-2.7B, Qwen2-7B, and Gemma-2-2B, although the magnitude varies across models (see Table~\ref{tab:dark_subspace}). For every model-SAE setting that passes the validity gate, the residual exceeds the reconstruction. A one-sided binomial sign test gives $p \approx 0.0078$, so the cross-architecture inversion is statistically supported, not merely qualitative. This supports qualitative generality of the ordering, not the same evidential status as the controlled Pythia experiments. Pythia-6.9B (mixed) and Pythia-12B each pool five SAE seeds, while the remaining settings rely on a single SAE seed. Settings with low reconstruction quality or abnormal
feature activation are retained as context rather than treated as clean generalization evidence. Per-setting confidence intervals, direction-of-effect
tests, and attempted replications are reported in Appendix~\ref{app:koc2_bootstrap} and Appendix~\ref{app:qwen2_pilot}.

Figure~\ref{fig:score_distributions} shows why the AUROC changes. In the Pythia-6.9B mixed-data setting, original activations separate members from non-
members. SAE reconstruction collapses much of that separation, and the reconstruction residual restores much of it. The effect is therefore visible not
only as a scalar AUROC reduction, but as a movement of the member and non-member score distributions.

\begin{figure}[!t]
\centering
\includegraphics[width=\linewidth]{figures/score_distributions_simple.pdf}
\caption{Per-text membership score distributions before and after SAE reconstruction for Pythia-6.9B, using the mixed-data SAE pooled across five
seeds. Original activations separate members from non-members (AUROC $=0.803$), SAE reconstruction collapses the separation (AUROC $=0.594$), and the
residual restores much of it (AUROC $=0.781$).}
\label{fig:score_distributions}
\end{figure}

\textbf{Variance structure of the residual readout.}
The variance across five Pythia-12B SAE initializations is consistent with the residual-over-reconstruction interpretation. As the SAE seed varies, the residual detector recovers membership signal at nearly constant strength, whereas the dictionary-projection detector is more variable. This asymmetry is evidence against a simple chance-miss explanation. Per-seed standard deviations and the algebraic identity relating drop variance to reconstruction variance are reported in Appendix~\ref{app:p12b_replication}. The magnitude of the original-vs-reconstruction drop is partition-sensitive under held-out estimation of $\dK$ (\S\ref{sec:results:controls}), whereas the inversion reported here is not.

\subsection{Decoder span, sparse codes, residual norm, and top features do not explain the residual signal}
\label{sec:results:adequacy}

We evaluate four alternative explanations for the residual membership signal, namely limited decoder span, preservation in sparse codes, dependence on residual norm alone, and concentration in a small set of high-activation features. Projecting activations onto the SAE decoder column space preserves the original membership detection AUROC on models with reliable reconstructions, while the orthogonal-complement projection falls near chance. The relevant direction is geometrically compatible with the decoder span. The loss of membership evidence happens in the sparse encode--decode path, not in the linear span alone.

A second possibility is that the sparse code preserves the signal even if the reconstruction does not. This also fails. A fold-matched logistic detector
trained on SAE codes $z$ does not recover the membership signal that the same detector extracts from the raw hidden state, with Holm--Bonferroni adequacy
failing on three of four valid models (Appendix~\ref{app:fresh_probe_v2}). These two tests narrow the interpretation. The issue is not that the SAE
decoder span is geometrically incapable, but that the learned sparse code and reconstruction do not preserve the membership evidence in a usable form.

\textbf{Norm and feature-subset tests.}
Residual magnitude explains part of the signal in some settings, but not all of it. L2-normalizing residual vectors preserves directional membership signal (Table~\ref{tab:l2_normalized}). Ablating top-$k$ SAE features ranked by point-biserial correlation further weakens reconstructed membership signal rather than recovering it. The feature sufficiency criterion gives the geometric version of the same result. Across all nine models, SAE classifier features span only a small fraction of the knowledge subspace $\SK$ (Table~\ref{tab:fsc_values}). On Pythia-6.9B, ablating all classifier features changes membership AUROC by less than a thousandth while collapsing verbatim extraction. Methods that select SAE features by training-set correlation~\citep{wang2025sspu,frikha2025privacyscalpel,muhamed2025dsg} therefore target the recall channel rather than the membership signal. Per-explanation verification protocols and detailed numerical evidence are in Appendix~\ref{app:sae_features} and Appendix~\ref{app:fsc}.

\subsection{Feature edits do not close the residual membership gap}
\label{sec:results:interventions}

None of the four alternatives explains the residual signal on its own. We next test whether feature ablations that suppress extraction also reduce membership detection.

\textbf{Feature ablation separates extraction from membership detection.}
Feature ablation shows that suppressing extractable recall is not the same as suppressing membership detection. Selected SAE features can suppress extractable recall while leaving membership detection within roughly 1 pp of baseline. On the four models tested at erasure ranks 10 and 20, removing leading
knowledge directions collapses verbatim extraction, but membership detection remains near ceiling (Appendix~\ref{app:dd_extraction}). This supports a
practical caution. An intervention that makes the model stop reproducing text need not remove the evidence used by a membership detector.

\textbf{Targeting the known memorization direction does not close the residual gap.}
We next test whether the residual gap disappears when the SAE is explicitly trained to reconstruct the known memorization direction $\dK$. Adding a
reconstruction penalty along $\dK$ restores the $\dK$-based score in the SAE reconstruction, showing that the targeted direction can be moved into the
reconstructed component. A general-purpose detector still recovers membership signal from the residual at a level comparable to the original activation
(Appendix~\ref{app:fresh_probe_v2}). A random-direction control does not produce the same recovery, indicating that the targeted effect is specific to $
\dK$.

The result suggests that $\dK$ is not the whole membership signal. Figure~\ref{fig:privacy_aware} summarizes this contrast on Pythia-6.9B. Standard SAEs
reduce score$_K$ after reconstruction, whereas direction-penalized SAEs (trained with an additional $\dK$-projection penalty) recover score$_K$ but leave membership signal detectable in the residual. Aligning the
SAE reconstruction with $\dK$ therefore does not show that membership information is confined to the targeted direction. Because this experiment uses
ground-truth $\dK$ and a single seed per condition on two models, it should not be read as a practical direction-penalized training protocol. Label-free proxy
directions and stronger replication would be needed before treating it as an intervention method.

\begin{figure}[!t]
\centering
\includegraphics[width=\linewidth]{figures/privacy_aware_comparison.pdf}
\caption{Membership detection before and after SAE reconstruction on Pythia-6.9B. Standard SAEs reduce score$_K$ after reconstruction, whereas direction-penalized SAEs (trained with an additional $\dK$-projection penalty) recover score$_K$ but leave membership signal detectable in the residual.}
\label{fig:privacy_aware}
\end{figure}

Across the feature-ablation and dissociation results, the SAE feature intervention target is not the activation direction that carries fine-tuning membership.

\subsection{Geometry and concentration of the residual signal}
\label{sec:results:controls}

We next characterize the direction and dimensionality of the residual signal, then test robustness to paraphrase, surface-form, and operating-point controls.

\textbf{Held-out estimation of $\dK$.}
Held-out estimation of $\dK$ does not support the original magnitude claim, but the qualitative residual $\geq$ original $>$ reconstruction ordering survives held-out estimation on both Pythia models (Appendix~\ref{app:heldout_dk_protocol}). The interpretation is therefore split. $\dK$ as a single direction is partition-sensitive in magnitude, but the ordering is not driven by evaluating and fitting that direction on the same documents.

\textbf{Residual-direction ablation is exploratory.}
We ask whether the residual membership signal is concentrated in a small number of directions. On Pythia-12B, ablating the top 10 residual principal
components reduces membership detection AUROC relative to random-rotation and matched-noise controls (Table~\ref{tab:kpc_kten_cells}). This suggests
partial concentration in the residual, but the intervention does not satisfy all falsification criteria, including the random-column-mask and same-SAE
consistency tests (Appendix~\ref{app:additional_controls}). We therefore treat it as evidence about residual geometry, not as causal localization of the
membership signal.

\textbf{Architecture and norm-direction structure.}
The norm-direction decomposition reveals a descriptive split between architecture families (Table~\ref{tab:norm_direction}, Appendix~\ref{app:per_model_bcd}). The GPT-family and LLaMA-family groups differ on whether membership separation comes from activation magnitude or from the channel direction, with implications for what SAE reconstruction preserves. The split is descriptive only because the LLaMA-family group contains only two models, and all LLaMA-family models in our sample postdate 2023. Full architecture scatter plots are in Figure~\ref{fig:arch_and_scaling}, and the per-layer concentration analysis is in Appendix~\ref{app:dimensionality_details}.

\subsection{Robustness to confound and operating point controls}
\label{sec:results:robustness}

We next test whether the main reconstruction-residual finding is robust to paraphrase, surface-form controls, and low-FPR operating points.

\textbf{Detector orientation under paraphrase.}
Under word-order-shuffled paraphrase, the residual detector orientation reverses sign on every model audited. A valid paraphrase audit therefore requires fixed-orientation AUROC. Per-model orientation diagnostics and the bidirectional-audit methodology are in Appendix~\ref{app:tpr_paraphrase}.

\textbf{Surface, access-setting, and operating-point controls.}
We evaluate the main findings against five families of confounds. For surface confounds, score$_K$ exceeds the activation norm baseline (Table~\ref{tab:norm_direction}), a length classifier and a shuffled-label null sit at chance, and the bag-of-words baseline rules out a token-level vocabulary account (Appendix~\ref{app:bow_baseline}). For corpus overlap, a corpus-disjoint dictionary on Pythia-6.9B argues against train-test SAE corpus overlap (Appendix~\ref{app:corpus_disjoint}). For within-family replication, Pythia-12B mixed-data tests whether the residual signal recurs in a separately trained setting (Appendix~\ref{app:p12b_replication}). For output-level scope, standard output-level detectors show no membership detection drop after SAE reconstruction, so the layer-local effect appears confined to the SAE layer (Appendix~\ref{app:standard_probes}). For operating point, TPR at 1\% and 0.1\% FPR is in Appendix~\ref{app:tpr_paraphrase}.